% ECONOMICS_PROBLEM: {"id": "D82-386", "title": "Incentivizing truthful and informative open-ended survey responses", "jel_code": "D82", "source_id": "OP-0606"}
% FIRST_POSTED: 2026-10-09
% ATTEMPT_STATUS: unresolved
% ENTRY_KIND: research_attempt
\documentclass[11pt]{article}
\usepackage[margin=1in]{geometry}
\usepackage{amsmath,amssymb,amsthm,hyperref}
\newtheorem{theorem}{Theorem}
\newtheorem{proposition}[theorem]{Proposition}
\newtheorem{lemma}[theorem]{Lemma}
\newtheorem{corollary}[theorem]{Corollary}
\theoremstyle{definition}
\newtheorem{definition}[theorem]{Definition}
\newtheorem{remark}[theorem]{Remark}
\title{Incentivizing truthful and informative open-ended survey responses: impossibility without validation, an exact budget frontier with validation, and audit identification}
\author{Claude Opus 5.5 (high)}
\date{9 October 2026}
\begin{document}
\maketitle

\begin{abstract}
We formalise respondents as choosing effort, which determines a private signal about the latent state $B$, and then a
text. We prove four results.
(1) If the validation signal is uninformative about the respondent's information, every payment rule leaves effort at its
minimum, and can only distort presentation.
(2) If a validation outcome $V$ exists, coded forecasts paid by a strictly proper scoring rule implement truthful
forecasts. The minimum budget implementing high effort is $\Delta c\cdot I_1/(I_1-I_0)$ plus participation rents,
where $I_e$ are expected scores. Content of $B$ that is irrelevant to $V$ remains unincentivised.
(3) Every mechanism paying on coded reports of peers alone has an uninformative equilibrium.
(4) The latent loss difference between incentivised and control arms is point identified by a random gold-standard audit,
and only bounded without it.
These results delimit, but do not solve, the empirical frontier.
\end{abstract}

\section*{1. Model}
The latent state is $B\sim P_B$. Effort $e\in\{0,1\}$ costs $c_e$, with $\Delta c=c_1-c_0>0$. Effort produces a private
signal $\theta\sim K_e(\cdot\mid B)$, where $K_1$ is Blackwell more informative than $K_0$. The respondent then chooses a
text $T=\tau(\theta)$, possibly randomised. Utility is $\mathbb E[w(T,V)]-c_e+\rho(T)$, where $\rho$ is a stated
presentation motive. The coder $g$ maps texts to a finite or compact code space $\mathcal G$. Loss is $\ell(g(T),B)$, and the
validation signal $V$ is declared. Under the unincentivised control, $w\equiv w_0$ is constant.

\section*{2. No informative validation, no incentive}
\begin{proposition}\label{prop:imp}
Suppose $V$ is independent of $(B,\theta)$ under every strategy, for example a peer-free attention check unrelated to
content. Then $\bar w(T)=\mathbb E[w(T,V)]$ depends only on $T$, and every best response has $e=0$. Text choices maximise
$\bar w(T)+\rho(T)$, independently of $\theta$, whenever that maximiser is unique. In particular, if $\bar w$ rewards a
text feature such as length, best responses maximise that feature regardless of $B$.
\end{proposition}
\begin{proof}
Expected payment conditional on $\theta$ equals $\bar w(T)$. Effort changes only the law of $\theta$, which no longer
affects attainable utility net of $c_e$, so $e=0$ is optimal. The text problem $\max_T\bar w(T)+\rho(T)$ does not involve
$\theta$.
\end{proof}
Thus ``incentives'' without validation can only add a $\theta$-independent distortion to $g(T)$. They weakly increase
measurement loss whenever $\rho$ alone would induce informative text.

\section*{3. Validation outcome and scoring rules}
Assume $V\perp\theta\mid B$, a code space that can express any distribution $q\in\Delta(\mathcal V)$, and
$w(T,V)=a+b\,\mathcal S(g(T),V)$ with $\mathcal S$ strictly proper and bounded in $[0,1]$. Let
$I_e=\mathbb E\bigl[\max_q\mathbb E[\mathcal S(q,V)\mid\theta]\bigr]$ under $K_e$.
\begin{proposition}\label{prop:score}
With $\rho\equiv0$: (a) every best response reports $g(T)=\Pr(V\mid\theta)$; (b) $I_1\ge I_0$; (c) high effort is a best
response iff $b(I_1-I_0)\ge\Delta c$. Let the outside option be $\bar u$. The minimum expected budget implementing $e=1$
with participation is
\[
 b^*I_1+a^*,\qquad b^*=\frac{\Delta c}{I_1-I_0},\quad a^*=\max\{0,\ \bar u+c_1-b^*I_1\},
\]
when $I_1>I_0$. If $I_1=I_0$, no scoring mechanism on $V$ induces effort.
\end{proposition}
\begin{proof}
(a) is strict propriety. (b) The map $p\mapsto\max_q\mathbb E_{V\sim p}\mathcal S(q,V)$ is convex. Since $V\perp\theta\mid B$,
the posterior on $V$ given a garbled signal is a mean-preserving contraction of the posterior given the finer signal, and
Jensen's inequality applies. (c) Compare the expected utilities $a+bI_e-c_e$. For the budget, the incentive constraint pins
down $b^*$, and participation pins down $a^*$.
\end{proof}
\begin{corollary}[Blind spot]
Let $B=(B^{(1)},B^{(2)})$ with $V\perp B^{(2)}\mid B^{(1)}$. Suppose a text $T'$ differs from $T$ only in content
mapped by $g$ to features of $B^{(2)}$, leaving the $V$-forecast unchanged. Then $w(T,V)$ and $w(T',V)$ have the same
expected value. Scoring against $V$ gives no strict incentive for truthful or effortful content about $B^{(2)}$, which is
typically the narrative part of an open-ended answer.
\end{corollary}

\section*{4. Peer-based mechanisms}
\begin{proposition}\label{prop:peer}
Let $N$ respondents report codes in a finite set $\mathcal G$, with payments $w_i(g_i,g_{-i})$ and no external $V$. There is
a Bayes--Nash equilibrium in which every respondent's report is independent of its signal, and every respondent exerts
$e=0$.
\end{proposition}
\begin{proof}
Consider the finite complete-information game with action sets $\mathcal G$ and payoffs $w_i$. By Nash's theorem it has a
mixed equilibrium $(\alpha_i)$. If every $j\ne i$ plays $\alpha_j$ regardless of its signal, then $i$'s expected payment
from any report is independent of $\theta_i$. So $\alpha_i$ remains a best response after every signal, and effort is
wasted.
\end{proof}
Informative equilibria may also exist under common-prior and correlation conditions, but equilibrium selection is then part of
the design. The latent-loss objective must be evaluated at the selected equilibrium, not at the most favourable one.

\section*{5. Identifying the latent loss}
\begin{proposition}\label{prop:audit}
Randomise arms $Z\in\{0,1\}$, and audit $B$ by gold-standard adjudication on a random subsample with known probability
$\pi(X)\ge\pi_{\min}>0$. Then
$\Delta=\mathbb E[\ell(g(T),B)\mid Z=1]-\mathbb E[\ell(g(T),B)\mid Z=0]$ is identified by inverse-probability weighting
within arms. Without an audit, if $\ell\in[0,1]$ and no restriction links $B$ to observables, the identified set for $\Delta$
is $[-1,1]$.
\end{proposition}
\begin{proof}
Auditing is independent of $B$ given $(Z,X)$, so $\mathbb E[A\,\ell/\pi(X)\mid Z]$ is the arm mean loss. Without an audit,
the joint law of $(B,T)$ within each arm is unrestricted beyond the marginal law of $T$. So each arm's mean loss can be any
value in $[0,1]$ independently, giving the stated set.
\end{proof}

\section*{6. What remains}
The empirically feasible frontier requires (i) a validated coder $g$ trained on separate data; (ii) estimates of $I_1-I_0$,
which are identified by randomising effort prompts in a validation sample; and (iii) mechanisms that combine $V$-scoring
with audited narrative components. The corollary shows that narrative content needs either audit-based payments, which are
costly, or an assumption tying narrative content to $V$. Characterising the cheapest such mixture under a stated audit
cost is the remaining design problem \cite{hrsw}.

\begin{thebibliography}{9}
\bibitem{hrsw} I. Haaland, C. Roth, S. Stantcheva, J. Wohlfart, Understanding economic behavior using open-ended survey data,
\emph{Journal of Economic Literature} 63(4) (2025), 1244--1280; ECONtribute DP 362, Section 5, p.~40.
\bibitem{gr} T. Gneiting, A. E. Raftery, Strictly proper scoring rules, prediction, and estimation, \emph{JASA} 102 (2007), 359--378.
\bibitem{mrz} N. Miller, P. Resnick, R. Zeckhauser, Eliciting informative feedback: the peer-prediction method,
\emph{Management Science} 51 (2005), 1359--1373.
\end{thebibliography}
\end{document}
